% Lean module: MonadFormal.Information.Testing.Sampling
\paragraph{Concentration under independent sampling.}
For a real statistic $f$ on the finite alphabet, put
$S_n(f)=\sum_{i=1}^n f(X_i)$, with the coordinates independently
distributed according to a probability law $R$.
\leandef{MonadFormal.Information.sampleSum}

\begin{lemma}[A quantitative typical set]
\label{lem:iid-typical-set}
Let $m=\mathbb E_Rf$ and $v=\operatorname{Var}_R(f)$.
Then $\mathbb E S_n(f)=nm$ and
$\operatorname{Var}(S_n(f))=nv$. For $n\ge1$ and $\varepsilon>0$,
\[
 R^{\otimes n}\{|S_n(f)-nm|\ge n\varepsilon\}
 \le\frac{v}{n\varepsilon^2}.
\]
In particular, if $2v\le n\varepsilon^2$, the event
$A_n=\{|S_n(f)-nm|<n\varepsilon\}$ has probability at least $1/2$.
\end{lemma}
\begin{proof}
The coordinate marginals give the mean by linearity of expectation.
Independence gives the variance sum; all moments are finite on the
finite alphabet. Chebyshev's inequality at deviation $n\varepsilon$
then gives $nv/(n\varepsilon)^2=v/(n\varepsilon^2)$.
Taking the complementary event proves the last assertion.
\leanproof{MonadFormal.Information.integral_sampleSum,MonadFormal.Information.variance_sampleSum,MonadFormal.Information.iid_deviation_le,MonadFormal.Information.iid_typical_ge_half}
\end{proof}
